% solver_result_contract.tex —— 求解路径、后验证与结果字段

\section{求解路径与结果证书}\label{sec:nr-solver}

\subsection{模型装配与二元变量}

矩阵以 Eigen 稀疏 triplet 装配为 \texttt{MIPModel}。只有自由 $\beta_e$ 和全部 $\gamma_g$ 被声明为
二元变量；固定 $\beta_e$ 保持连续类型但满足 $lb=ub\in\{0,1\}$，以减少 B\&C 分支规模。
故障边、不可操作边、排除 overlay 和同分量无效 tie 均通过变量边界固定。

\fld{enable\_pf=false} 时只创建 $F,F_g,\beta,\gamma$；此时
\fld{enable\_voltage}、\fld{enable\_thermal} 和 \fld{loss\_aware} 不创建相应变量或约束。
\fld{model\_scope} 分别固定为：
\begin{equation}
\begin{cases}
\texttt{hybrid-acdc-topology-lindistflow},&\fld{enable\_pf}=true,\\
\texttt{hybrid-acdc-topology-connectivity},&\fld{enable\_pf}=false.
\end{cases}
\end{equation}

\subsection{图启发式}

当 \fld{skip\_heuristic=false} 且不是分域模式时，核心先从主根沿当前健康闭合边做 BFS。若所有母线
已可达，则取 BFS 树；若有不可达节点，则从一条连接“已达/未达”两侧的初始断开自由边中选择当前
线性系数最小者，再扩展一次 BFS。启发式向量设置：
\begin{itemize}
  \item participating 边只闭合 BFS parent edge，overlay 保留原状态；
  \item 主根 $\gamma=1$，虚拟流由子树大小生成；
  \item 其余变量初始为零。
\end{itemize}
只有该向量同时满足全部等式、不等式和变量界（最大违反 $<10^{-6}$）时才被接受。PF 开启且存在
非零负荷时，零 $P/Q/P_g/s$ 往往使启发式失败，从而进入 MILP。

成功的启发式返回
\begin{center}
\texttt{solver\_backend=graph-heuristic},\quad
\texttt{solver\_status=feasible incumbent},\quad
\texttt{solver\_mip\_gap=inf},\quad
\texttt{proven\_optimal=false}.
\end{center}
它没有与 MILP 下界比较，不能称为最优 ONR。

\subsection{实际后端分派}

\begin{center}\small
\begin{tabular}{@{}L{2.5cm}L{6.0cm}L{6.0cm}@{}}
\toprule
\fld{solver} & 首选路径 & 失败行为\\
\midrule
\texttt{native} & \texttt{solve\_milp\_bc}，MIR、5 轮根割、伪成本分支 & 直接失败返回\\
\texttt{scip} & \texttt{ScipAdapter} & 不可用或失败则返回\\
\texttt{highs} & \texttt{StrictHighsBranchAndCutAdapter} & 再尝试 SCIP；不尝试 Native\\
\texttt{auto} & HiGHS & 再尝试 SCIP；不尝试 Native\\
其他字符串 & 与 \texttt{auto} 相同 & HiGHS 后 SCIP\\
\bottomrule
\end{tabular}
\end{center}

这一分派以 \srcpath{topology\_reconfiguration.cpp:run\_topology\_reconfiguration} 为准。
公共头注释中 \texttt{auto (HiGHS->native)} 与实现不符；当前实际没有 auto-to-native fallback。
外部后端成功时 \fld{bc\_stats} 保持默认值，诊断应读取 \fld{solver\_backend}、
\fld{solver\_status} 和 \fld{solver\_mip\_gap}。

\subsection{成功、可行与最优判定}

外部或 Native 返回 \texttt{stats.success=true} 且解向量长度不少于变量数，才进入统一后验证。
\fld{proven\_optimal} 还要求：
\begin{enumerate}
  \item 规范化状态字符串属于实现识别的最优状态集合；
  \item MIP gap 有限；
  \item $gap\le\fld{mip\_gap}+10^{-9}$。
\end{enumerate}
识别集合包括 \texttt{optimal}、\texttt{highs optimal}、
\texttt{optimal solution found}、\texttt{optimality gap reached} 等固定字符串。一个后端若返回语义等价
但未列入的状态，仍会得到可行结果，但不会标记已证明最优。

\subsection{统一后验可行性检查}

所有启发式/外部/Native 解必须通过：
\begin{align}
\|A_{eq}x-b_{eq}\|_\infty&\le10^{-6},\\
\|(Ax-b)_+\|_\infty&\le10^{-6},\\
\operatorname{dist}(x_j,\{0,1\})&\le10^{-4}\quad(j\in\mathcal B),\\
l_i-10^{-6}\le x_i&\le u_i+10^{-6}.
\end{align}
任何一项失败都把 \fld{feasible} 改为 false 并返回；残差数值只在 verbose 日志中出现，没有存入
\texttt{TopoReconfResult}。通过后 \fld{feasible=true}，且
\fld{optimal=proven\_optimal}。

\subsection{目标值与分解}

模型内开关项对初始闭合边使用 $-\lambda_{sw}c_e\beta_e$，省略常数
$\lambda_{sw}c_e$；初始断开边使用 $+\lambda_{sw}c_e\beta_e$。因此
\fld{milp\_objective}$=c^Tx$ 可以为负，且不是四个正向业务成本的直接总和。

\texttt{ObjBreakdown} 采用更易读的正向动作计费：
\begin{align}
J_{switch}^{report}&=\lambda_{sw}\sum_e c_e[\beta_e\ne\beta_e^0],\\
J_{loss}^{report}&=\lambda_{loss}\sum_{e\in E_{line}}|r_e|
\begin{cases}t_e,&\fld{loss\_aware}\land\fld{enable\_pf},\\\beta_e,&\text{其他},\end{cases}\\
J_{shed}^{report}&=\lambda_{shed}\left(\sum_i s_i^P+\sum_{i\in AC}s_i^Q\right),\\
J_{island}^{report}&=\lambda_{island}\sum_{g\ne g_0}\gamma_g.
\end{align}
因为 $J_{switch}^{report}$ 恢复了省略常数，通常有
\begin{equation}
\fld{milp\_objective}\ne
J_{switch}^{report}+J_{loss}^{report}+J_{shed}^{report}+J_{island}^{report}.
\end{equation}
调用方不得用四项求和反向校验原始目标，除非同时恢复所有常数。

\subsection{损耗字段的三种不同口径}

\begin{center}\small
\begin{tabular}{@{}L{4.0cm}L{5.1cm}L{5.4cm}@{}}
\toprule
字段/数值 & 计算 & 含义\\
\midrule
\fld{obj\_terms.loss} & $\lambda_{loss}\sum |r|t$ 或 $\lambda_{loss}\sum|r|\beta$ & 加权目标贡献，无物理 MW 承诺\\
\fld{reconf\_loss\_mw} & $S_B\sum_{AC/DC\ closed}r_e$ & 假设每条闭边 1 pu 电流的拓扑代理\\
HTTP \texttt{reconfig\_loss\_mw} & 完整 PF 支路 $P_{from}+P_{to}$ 的正损耗和 & PF 收敛时的后验证损耗\\
\bottomrule
\end{tabular}
\end{center}

核心 \fld{base\_loss\_mw}、\fld{loss\_reduction\_mw}、\fld{loss\_reduction\_pct} 当前没有赋值
路径，保持默认 0；HTTP 使用局部变量另行计算并写入 JSON。\texttt{ONRResult::estimated\_loss\_mw}
映射自核心 \fld{reconf\_loss\_mw}，仍是拓扑代理。

\subsection{完整字段赋值真值表}

\begin{longtable}{@{}L{4.2cm}L{4.0cm}L{6.4cm}@{}}
\toprule
字段 & 核心赋值时机 & 解释/限制\\
\midrule
\endfirsthead
\toprule
字段 & 核心赋值时机 & 解释/限制\\
\midrule
\endhead
\fld{feasible} & 后验线性检查通过 & 不是 PF/OPF 可行\\
\fld{optimal}/\fld{proven\_optimal} & 状态+有限 gap 通过 & 启发式永远 false\\
结构化 open/closed/switched & 可行解提取 & 混合系统权威索引空间\\
旧式整数向量 & 同时填充 & 可能跨域碰撞\\
\fld{switch\_operations} & 状态发生变化 & 设备或分支稳定 ID\\
\fld{ordered\_switch\_actions} & 动作回投影后 & 静态操作顺序，不是时域仿真\\
\fld{total\_shed\_mw/mvar} & PF 变量存在时 & $s^P/s^Q$ 求和；无节点向量\\
\fld{milp\_objective} & $c^Tx$ & 含权重和省略常数\\
\fld{obj\_terms} & 可行解后处理 & switching 项与原目标常数口径不同\\
\fld{reconf\_loss\_mw} & 可行解后处理 & 名义电流代理\\
\fld{base\_loss\_mw} 等 & 不赋值 & 核心始终默认 0\\
\fld{solve\_time\_s} & 正常/大多数失败返回 & 早期空系统/无源返回仍为 0\\
\fld{solver\_backend/status/gap} & 求解路径 & 早期无源只有 status\\
\fld{bc\_stats} & 仅 Native 路径 & 外部后端保持默认\\
\fld{model\_scope} & 入口开始 & 即使后续不可行仍存在\\
\fld{post\_action.../post\_power.../full\_hybrid.../executable} & 不赋值 & 仅 HTTP 层修改副本\\
\bottomrule
\end{longtable}

\subsection{ValidityFlags 的准确读法}

核心入口在求解前即设置部分标志：\fld{radial\_topology\_enforced}、
\fld{device\_capability\_constraints\_enforced}、\fld{protection\_interlocks\_enforced} 为 true；
\fld{ac\_lindistflow\_enforced} 等于 \fld{enable\_pf}；DC/VSC 标志按资产存在设置。
这些标志描述“模型构造意图”，并不以 \fld{feasible} 为前提。

特别地，\fld{allow\_dc\_mesh=true} 时 DC 边不受树边数约束，但
\fld{radial\_topology\_enforced} 仍为 true。调用方必须联合读取选项/运行时拓扑结果，不能把该单个
布尔值解释为 AC、DC、VSC 全域都已证明径向。

\subsection{失败模式}

\begin{itemize}
  \item 空母线或无 AC/DC 线路：返回默认不可行结果；只有 VSC 而无线段也在此边界内；
  \item 无可用根源：返回不可行并给出明确 \fld{solver\_status}；
  \item 后端无 incumbent：返回不可行并保留后端失败摘要；
  \item 后验残差、整数性或边界失败：返回不可行，详细残差只在 verbose 日志；
  \item 无效故障/候选稳定 ID：混合核心多数静默忽略；历史 ONR 旧死代码中的 throw 不可达；
  \item 未知 \fld{solver}：落入 auto 行为，不拒绝输入。
\end{itemize}
